\PID \label{z:mnk_1}
Познато је да физичка величина $y$ зависи од физичке величине $x$ према линеарној 
зависности $y = kx$. Ако је познато $n$ резултата добијених мерењима $(x_1, y_1)$, 
$(x_2, y_2)$, \ldots, 
$(x_n, y_n)$, одредити оптималну вредност параметра $k$, тако да је сума квадрата одступања
од измерених вредности минимална (\textit{Метод најмањих квадрата}).

\RESENJE

Одступање $i$-тог мерења се може записати у облику $e_i = kx_i - y_i$. Сума квадрата свих одступања 
$\upepsilon = \sum_{i = 1}^n e_i^2$ је 
онда 
\begin{equation}
    \upepsilon = \sum_{i = 1}^n (kx_i - y_i)^2  
    = \sum_{i = 1}^n (k^2 x_i^2 + y_i^2 - 2 k x_i y_i)
    = k^2 \sum_{i = 1}^n x_i^2 - 2 k \sum_{i = 1}^n (x_i y_i) + \sum_{i = 1}^n y_i^2
\end{equation}
Добијена сума одступања представља квадратну зависност од параметра $k$, која због позитивног најстаријег коефицијента
представља конвексну параболу, која има своју минималну вредност. Минимална вредност се може потражити и одређивањем нуле извода, па је
коначно 
\begin{equation}
    \dfrac{\partial  \upepsilon  }{\partial k} =
    2 k \sum_{i = 1}^n x_i^2 - 2 \sum_{i = 1}^n (x_i y_i) = 0 
    \Rightarrow
    k = \dfrac{\displaystyle \sum_{i = 1}^n (x_i y_i) }{\displaystyle \sum_{i = 1}^n x_i^2} \label{eq:mnk_1:rezultat}
\end{equation}
